\chapter{The test suite: seams, fixtures, gates, and how a timing test must be written}\label{ch:tests}

The suite is 1591 tests on Linux with the offscreen platform, plus the \texttt{examples/} harness that
CI also runs. As this manual was written it ran green in 336.28 s: 1586 passed, 5 skipped (the two
build-gate tests, the two container-gate tests, and the opt-in review-screenshot test). It is long
(about five and a half minutes on this machine) and it is the only place most of the rules in this
manual are actually proven, so a change to it is as consequential as a change to the code.

\section{Fixtures that every test inherits}

\texttt{conftest.py} is not boilerplate; each fixture in it exists because of a failure.

\begin{itemize}
  \item \textbf{One \texttt{QApplication} for the session.} This Qt 6.11 build segfaults at interpreter
    shutdown when a live \texttt{QApplication} is finalised alongside any \texttt{QCoreApplication}.
    Sharing one application means only a lone application is ever finalised. The binding's
    \texttt{QApplication} does not accept a \texttt{platform=} keyword, so the offscreen platform comes
    from \texttt{QT\_QPA\_PLATFORM}.
  \item \textbf{Widget trees and workers are torn down per test.} A \texttt{StagedRunScreen} keeps a live
    worker and a one-second timer for its whole lifetime; if it survives into the next test, its queued
    signals fire into that test's widgets and corrupt their visibility --- the \texttt{run\_view} flake.
    \texttt{\_stop\_workers} and \texttt{\_\allowbreak{}destroy\_\allowbreak{}widget\_\allowbreak{}tree} run before every test, and the widget
    tree is destroyed while Python is still alive.
  \item \textbf{Every started engine is stopped after every test.} \texttt{\_\allowbreak{}stop\_\allowbreak{}running\_\allowbreak{}engines}
    wraps \texttt{RunEngine.start} weakly, then aborts and joins anything still running. If it cannot, the
    test fails with ``a RunEngine worker outlived its test'' and the description of the stuck run
    (section below).
  \item \textbf{Files are closed, strictly.} On Linux an open file can be deleted; on Windows it cannot.
    The \texttt{\_release\_files} fixtures close every \texttt{RunStore}, \texttt{HeadlessRunSession} and
    \texttt{AppLog} a test opened, then assert that nothing under pytest's temp directory is still open.
    A handle that only Windows would complain about is caught here, on Linux.
  \item \textbf{The data directory is redirected.} Config defaults are written as
    \texttt{<data>/\dots}; the autouse fixture points \texttt{ate.config.\_\allowbreak{}platform\_\allowbreak{}data\_\allowbreak{}dir} at
    pytest's basetemp so no test writes to \texttt{/var/lib} or \texttt{C:\textbackslash ProgramData}.
  \item \textbf{Privilege is pinned.} \texttt{ATE\_ELEVATED=0} for the session, session-scoped on purpose:
    a module-scoped fixture that starts a run runs before any per-test fixture, so a function-scoped pin
    would arrive too late. A test that wants to be privileged sets \texttt{ATE\_ELEVATED=1}.
  \item \textbf{The PDF helper is a stub by default.} \texttt{\_fake\_pdf\_helper} replaces
    \texttt{ate.report.pdf\_\allowbreak{}job.pdf\_\allowbreak{}helper\_\allowbreak{}command} with a process that writes a 13-byte stub, except
    in the four tests that must prove a PDF has text. A test needing other helper behaviour patches it
    itself.
  \item \textbf{Qt's fatal messages are not swallowed.} \texttt{report\_qt\_fatal} keeps Qt's fatal
    message from dying inside pytest-qt's capture.
\end{itemize}

\texttt{tests/\_fixtures.py} holds the seams a test can use: \texttt{LoopbackPeer} and
\texttt{UdpLoopbackPeer} for Device tests (a TCP/UDP peer that records what it received),
\texttt{scratch\_directory} for temporary directories that conftest removes,
\texttt{executable\_stub} for a stand-in binary, \texttt{fake\_pdf\_helper} and
\texttt{recording\_\allowbreak{}pdf\_\allowbreak{}helper}, \texttt{wait\_\allowbreak{}for\_\allowbreak{}disconnect}, \texttt{watch\_store\_writes} (which
records every store call and its duration), \texttt{describe\_\allowbreak{}stuck\_\allowbreak{}runs},
\texttt{wait\_\allowbreak{}for\_\allowbreak{}aborted\_\allowbreak{}run}, \texttt{SlowAbortWarning} and
\texttt{assert\_\allowbreak{}no\_\allowbreak{}swallowed\_\allowbreak{}exception}.

\section{Two gates}

\begin{lstlisting}[language=ATEShell]
ATE_BUILD_TESTS=1 pytest -m build -q      # the real PyInstaller bundle, and the .deb built from it
ATE_DOCKER_TESTS=1 pytest -m build -q     # the debian:stable container smokes
\end{lstlisting}

Both are off by default, and both are marked \texttt{build}. The marker description in
\texttt{pyproject.toml} is the documented contract: a test marked \texttt{build} is skipped unless
\texttt{ATE\_BUILD\_TESTS=1}; the container tests skip unless \texttt{ATE\_DOCKER\_TESTS=1} and Docker is
present. \texttt{tests/\allowbreak{}test\_\allowbreak{}release\_\allowbreak{}ci\_\allowbreak28c.py} asserts that the gate is written the way it says, so a
container test cannot quietly start running on every CI job.

The five skips in a default run are therefore
\texttt{test\_\allowbreak{}pyinstaller\_\allowbreak{}build\_\allowbreak27a.py}'s frozen bundle,
\texttt{test\_\allowbreak{}linux\_\allowbreak{}packages\_\allowbreak28b.py}'s \texttt{.deb} built from it,
\texttt{test\_\allowbreak{}manuals.py}'s container permissions smoke,
\texttt{test\_\allowbreak{}release\_\allowbreak{}ci\_\allowbreak28c.py}'s three containers, and
\texttt{test\_\allowbreak{}finished\_\allowbreak{}run\_\allowbreak{}layout\_\allowbreak30b.py}'s review screenshots, which need
\texttt{ATE\_\allowbreak{}30B\_\allowbreak{}SHOTS}. Both gates were run for this manual; the numbers are in
chapter~\ref{ch:packaging}.

\section{Timing tests, and the rule that came out of 31e}

Packet 31e was two tests that failed intermittently on Windows CI and passed on Linux. Reading the notes
and the code together, the rule that emerged is worth applying to every new timing test:

\begin{enumerate}
  \item \textbf{Start the clock on an event, not on a sleep.} The old abort test started timing after
    \texttt{time.sleep(0.1)} with no proof the action had begun, so any stall of the worker around that
    point was counted as abort latency. \texttt{\_drive\_abort} now starts on the worker being inside
    \texttt{ATE.sleep}'s wait, or on \texttt{input\_asked} for a blocked input, and stops at
    \texttt{FINISHED}, before any read.
  \item \textbf{Separate the requirement from the expectation.} Story 83 says an abort should stop within
    0.5 s and must stop within 10 s. A test fails past 10 s and \textbf{warns} past 0.5 s, so a slow
    runner stays visible in CI's summary instead of being a false failure.
    \texttt{wait\_\allowbreak{}for\_\allowbreak{}aborted\_\allowbreak{}run} implements both, and prints the worker's timeline.
  \item \textbf{A stamp means what the code did, not what the test expected.} \texttt{record\_log} takes
    \texttt{time.monotonic()} when the line is recorded; the drain maps it to wall-clock later. A test
    that compares \texttt{at} with a deadline measures the scheduler, not the product. The fixed tests
    bracket each \texttt{record\_log} call on the worker and require \texttt{at} to fall inside the
    bracket.
  \item \textbf{Show that the bound can fail.} Every bound in these tests has a test that plants a stall
    and proves the test catches it: a 1 s injected stall gives ``abort took 1.011 s''; a busy loop of
    \texttt{1e10} iterations fails the 10 s requirement; a stamp taken at drain time fails the bracket.
    A bound nobody has seen violated is a bound nobody has tested.
  \item \textbf{Do not widen a timeout.} Packet 32b's checklist says it plainly: the 10 s wait stays and
    can still fail, and when it does it now says what the worker was doing.
  \item \textbf{Run under load.} The convention the packets use is ten repetitions of the changed files
    with six busy loops on three cores. Several flakes only appear under contention, and two shared-helper
    races (the \texttt{LoopbackPeer} that stopped reading at \texttt{stop()} and the VISA tests that read
    \texttt{peer.received} before the peer had seen the connection end) were found that way.
\end{enumerate}

\section{What \texttt{describe\_\allowbreak{}stuck\_\allowbreak{}runs} gives you}

When a run does not end, the failure message includes a report section: every live worker and session,
what is recorded in the store, the store's write timings with any write still in flight, and all thread
stacks from \texttt{faulthandler}. \texttt{pytest\_\allowbreak{}runtest\_\allowbreak{}makereport} adds it to any test that failed
with a worker running, whichever wait helper it used. Its purpose is to answer the question packet 31e
could not: whether a stall is the disk (long or in-flight writes) or something else (normal writes, with
the worker's stack elsewhere).

\section{Where a new test goes}

Name it after the packet that added it, in the style the suite already uses
(\texttt{tests/\allowbreak{}test\_\allowbreak{}device\_\allowbreak{}abort\_\allowbreak32b.py}). Put shared helpers in \texttt{tests/\_fixtures.py}, not in a
test file, so the next test does not copy a helper that has a known race. If your test starts a Qt process,
use \texttt{assert\_\allowbreak{}no\_\allowbreak{}swallowed\_\allowbreak{}exception}: exit 0 no longer means nothing went wrong.

\cando{write a timing test that can fail, and pick the right gate for a slow one.}
